\subsection*{3.2. Determinación de $\Psi^1_{L_\bullet}(A)$}

Consideremos un objeto $(X,\alpha)$ (3.1.3) de
$\underline\Psi_{L_\bullet}(A)$. Consideremos el complejo de
cocadenas de longitud uno
\begin{equation}\tag{3.2.1}
 [X\longrightarrow L_0]:\quad
 0\longrightarrow X\longrightarrow L_0\longrightarrow0\longrightarrow
 \cdots
\end{equation}
definido por el morfismo $j:X\longrightarrow L_0$, con $X$ como su
componente de grado cero. Es una resolución de $A$ y, por tanto, puede
identificarse con $A$ como objeto de la categoría derivada
$D(\underline C)$ [5]. El diagrama (3.1.3) da entonces un morfismo de
complejos de grado uno
\begin{equation}\tag{3.2.2}
 L_\bullet\longrightarrow[X\longrightarrow L_0],
\end{equation}
explicitado por el siguiente diagrama de cuadrados anticonmutativos:
\begin{equation}\tag{3.2.3}
\begin{tikzcd}[column sep=large,row sep=huge]
 \cdots & L_2\arrow[r,"d_1"]
 & L_1\arrow[r,"d_0"]\arrow[d,"u^{-1}=-\alpha"']
 & L_0\arrow[r]\arrow[d,"u^0=\mathrm{id}"'] & 0 & \cdots\\
 \cdots & 0\arrow[r] & X\arrow[r,"j"] & L_0\arrow[r] & 0 & \cdots
\end{tikzcd}\quad .
\end{equation}
(Utilizamos la convención usual que asocia al complejo de cadenas
$L_\bullet$ un complejo cuyo diferencial tiene grado $1$ y cuya
componente de grado $i$ es $L_{-i}$.) Obtenemos así un morfismo de
grado uno en la categoría derivada
\begin{equation}\tag{3.2.3}
 c(X,\alpha):L_\bullet\longrightarrow A
 \qquad
 \bigl(\text{un morfismo de grado uno en }D(\underline C)\bigr),
\end{equation}
que manifiestamente depende solo de la clase de isomorfismo de
$(X,\alpha)$. Hemos definido con ello una aplicación canónica
\begin{equation}\tag{3.2.4}
 \Psi^1_{L_\bullet}(A)\longrightarrow\Ext^1(L_\bullet,A)
 \overset{\mathrm{def}}{=}
 \Hom_{D(\underline C)}(L_\bullet,A[1]).
\end{equation}

\underline{Teorema} 3.2.5. \underline{La aplicación} (3.2.4)
\underline{es un isomorfismo de grupos, funtorial en el objeto}
$A$ \underline{de} $\underline C$.

Dejamos al lector la verificación de que (3.2.4) es aditiva y
funtorial en $A$, y nos limitaremos a demostrar que es biyectiva. Para
la inyectividad, mostremos primero que si el morfismo de complejos $u$
de (3.2.3) es cero en $D(\underline C)$, entonces ya es cero en la
categoría de homotopía $K(\underline C)$; es decir, existe un morfismo
\[
 s:L_0\longrightarrow X
\]
tal que
\begin{equation}\tag{3.2.5.1}
 js=u^0=\mathrm{id}_{L_0},
 \qquad sd_0=-u^{-1}=\alpha.
\end{equation}
En efecto, la hipótesis dice que la conclusión se vuelve cierta
después de componer el morfismo de complejos considerado con un
cuasi-isomorfismo adecuado $v$
\begin{equation}\tag{*}
\begin{tikzcd}[column sep=large,row sep=huge]
 0\arrow[r] & X\arrow[r,"j"]\arrow[d]
 & L_0\arrow[r]\arrow[d] & 0\\
 0\arrow[r] & C^0\arrow[r,"j'"] & C^1\arrow[r] & C^2\arrow[r] & \cdots
\end{tikzcd}\quad ,
\end{equation}
(y, sustituyendo si es necesario $C^i$ por $0$ para $i\geq2$ y
$C^1$ por
$\operatorname{Im}(C^0\longrightarrow C^1)
=\operatorname{Ker}(C^1\longrightarrow C^2)$, podemos suponer que
$C^i=0$ para $i\geq2$). Pero el morfismo (*) induce la identidad sobre
$A$ y, por tanto, identifica $X$ con el producto fibrado de $L_0$ y
$C^0$ sobre $C^1$. En consecuencia, un morfismo
$s':L_0\longrightarrow C^0$ que define una homotopía $vu\sim0$, y que
por tanto satisface $j's'=v^1u^0$, se factoriza a través de un morfismo
$s:L_0\longrightarrow X$ que satisface
$js=u^0=\mathrm{id}_{L_0}$. Así, $s$ escinde la extensión $X$ de
$L_0$ por $A$, de modo que el complejo $X\longrightarrow L_0$ es
isomorfo \underline{en} $K(\underline C)$ al complejo $K^\bullet$
concentrado en $A$. Para este último complejo es claro que
\[
 \Hom_{K(\underline C)}(L_\bullet,K^\bullet)
 \longrightarrow
 \Hom_{D(\underline C)}(L_\bullet,K^\bullet)
\]
es un isomorfismo, lo cual demuestra la afirmación.

Ahora, la elección de un morfismo
$\sisi{j}{s}:L_0\longrightarrow X$ que satisface (3.2.5.1) equivale
precisamente a una trivialización de la extensión $X$ de $L_0$ por
$A$ compatible con la trivialización parcial $\alpha$. Esto demuestra
la inyectividad de (3.2.4).

Para la sobreyectividad, partamos de un elemento del miembro derecho.
Por definición, puede describirse mediante un morfismo de complejos de
grado uno
\begin{equation}\tag{3.2.5.2}
\begin{tikzcd}[column sep=large,row sep=huge]
 \ldots & L_2\arrow[r,"d_1"]
 & L_1\arrow[r,"d_0"]\arrow[d,"u^{-1}"']
 & L_0\arrow[r]\arrow[d,"u^0"'] & 0\\
 & 0\arrow[r] & C^0\arrow[r] & C^1\arrow[r] & \ldots
\end{tikzcd}
\end{equation}
con cuadrados anticonmutativos, donde $C^\bullet$ es una resolución de
$A$ y, como antes, podemos suponer que $C^i=0$ para $i\geq2$. Así,
$C^0$ es una extensión de $C^1$ por $A$, y podemos formar la extensión
$X$ obtenida por imagen inversa a lo largo de
$L_0\longrightarrow C^1$. La anticonmutatividad en (3.2.5.2) muestra
que $-u^{-1}$ se factoriza a través de un morfismo
$\alpha:L_1\longrightarrow X$ que es una escisión parcial, es decir,
$j\alpha=d_0$ (con $j$ como en (3.1.2)). Además, la relación
$u^{-1}d_1=0$ de (3.2.5.2) implica $\alpha d_1=0$. Por tanto,
$(X,\alpha)$ es efectivamente un objeto de
$\underline\Psi_{L_\bullet}(A)$. Por construcción, el morfismo
(3.2.5.2) es la composición del morfismo (3.2.3) asociado a
$(X,\alpha)$ y del morfismo
\[
 [X\longrightarrow L_0]\longrightarrow C^\bullet
 =[C^0\longrightarrow C^1]
\]
cuyas componentes son $(f,u^0)$, donde $f:X\longrightarrow C^0$ es el
morfismo canónico obtenido de la definición de $X$ como imagen inversa
de la extensión $C^0$.

El morfismo precedente es un morfismo de resoluciones de $A$ e induce
la identidad sobre $A$. Se sigue, por la definición
$\Ext^1=\Hom^1_{D(\underline C)}$, que el elemento dado de
$\Ext^1(L_\bullet,A)$ descrito por (3.2.5.2) es precisamente
$c(X,\alpha)$. Esto demuestra la sobreyectividad de (3.2.4) y concluye
la demostración de 3.2.5.
